\documentclass{article}
\usepackage[T1]{fontenc}
\usepackage{lmodern}
\usepackage{graphicx}
\usepackage{amsmath,amssymb}
\usepackage[a4paper,margin=2cm]{geometry}
\usepackage[absolute,overlay]{textpos}
\setlength{\TPHorizModule}{1mm}
\setlength{\TPVertModule}{1mm}
\usepackage[indonesian]{babel}
\usepackage{hyperref}
\setlength{\emergencystretch}{2em}
% unit: Halaman 90

\begin{document}

% === Halaman 90 (llm) ===

\section*{Kriptanalisis Affine Cipher}

\begin{itemize}
    \item Affine Cipher:
    \[ \text{Enkripsi: } C \equiv mP + b \pmod{n} \]
    \[ \text{Dekripsi: } P \equiv m^{-1} (C - b) \pmod{n} \]
    
    \item Affine cipher mudah diserang dengan \textit{known-plaintext attack} (sebagian plainteks dari pesan diketahui)
    
    \item Misalkan kriptanalis mempunyai dua buah plainteks, $P_1$ dan $P_2$, yang berkoresponden dengan cipherteks $C_1$ dan $C_2$,
    
    \item maka $m$ dan $b$ mudah dihitung dari buah kekongruenan simultan berikut ini:
    \begin{align*}
        C_1 &\equiv mP_1 + b \pmod{n} \\
        C_2 &\equiv mP_2 + b \pmod{n}
    \end{align*}
\end{itemize}

\end{document}
